% Appendix: reproducibility and release.
\section{Reproducibility and Release}
\label{app:reproducibility}

\paragraph{Freeze.} Before each block of runs, the source code, configuration, and scripts
are hashed into a freeze record that every runner checks before making a call. Any change
requires a new record and an account of which results it invalidates. Changes affecting a
reported number are logged with the previous and revised values and the cells rebuilt or
relabeled, including the two label audits in Section~\ref{sec:setup}.
Sampling seeds are fixed (assertion samples: seed 0).

\paragraph{Provenance.} A shared client caches every model response under a hash of its
provider, model, prompt, and parameters; each run writes a manifest and call log. Every number
in the paper maps to a measurement-ledger row identifying the producing command and run.
All generative calls of the systems and of the treatment used \res{GEN-model}, between 28 September and
\ifnwdecon 7\else 5\fi{} October 2026. The provider's model name exposes no finer revision.
Cached responses reproduce the released artifacts exactly; fresh calls need not produce
identical generations.\ifnwdecon{} On the treated text, \res{DECON-recovered} page of the
full-text diagnostic returned a list with an incorrect closing bracket in three attempts.
The runner corrected the bracket in the recorded first response and applied its normal
validation and rendering; no page was regenerated.\fi{} Appendix~\ref{app:prompts} presents
one logged exchange for each generative stage, including the system prompt, input, and
output. Claude Opus~5.5 finalized all reported AI annotation labels after evidence review.
The prompt and response defects documented in Appendix~\ref{app:prompts} remain part of
the archived run configuration. The reported results are not measurements of a repaired
prompt builder or citation validator.

\paragraph{Commands.} With the corpus parsed, the reported artifacts are produced by:
\begin{itemize}[leftmargin=*,itemsep=1pt,topsep=2pt]
  \raggedright
  \item the reported \sys build of all five cutoffs:\\ \texttt{\small run\_series.py anne@v2 -{}-upto 5}
    (the first build is the series \texttt{\small anne});
  \item a comparison cell at cutoff $t$:\\ \texttt{\small run\_baseline.py anne \{X1|B1|B2|B4|X2\} -{}-upto t},
    with \texttt{\small -{}-chunk-words 150 -{}-budget-tokens 20000} for B2 and
    \texttt{\small -{}-graph anne@v2 -{}-budget-tokens 20000} for B4, and
    \texttt{\small b3\_lightrag.py anne -{}-upto t -{}-mode mix};
  \item the assertion inventory and sample of a cell:\\ \texttt{\small assertion\_inventory.py S -{}-upto t} and
    \texttt{\small assertion\_precision.py S -{}-upto t -{}-n 40 -{}-population all};
  \item disclosure candidates:\\ \texttt{\small m4\_tree.py S -{}-upto t};
  \item the lexical screen:\\ \texttt{\small artifact\_exposure.py S};
  \item the pair rule and the timeline citations applied to a cell's labels:\\ \texttt{\small apply\_pair\_rule.py S B5};
  \item table values:\\ \texttt{\small paper\_numbers.py}; their spread over cutoffs:\\ \texttt{\small summary\_stats.py};
  \item cost of the cells on disk:\\ \texttt{\small run\_cost.py S -{}-final}.
\end{itemize}

\paragraph{Release.} An anonymized archive provides the submission materials; the repository
address is \anonrepo. It contains the code, configuration, dependency specification and freeze records;
prompt templates and logged exchange excerpts; gold facts with evidence identifiers and reveal volumes;
and AI evaluation label files
(recall, annotated precision samples, citation sufficiency, disclosure verdicts with the
sentences assessed, and audit corrections). It also includes per-run cost summaries, every
system's generated wikis at every cutoff, both builds of \sys, and the human-review package
(\res{ANN-n-total} items with source context: every gold fact and samples of the recall,
precision and citation labels), plus the two reviewers' label files, accompanying records,
provenance manifest, aligned decision audit and recomputed agreement. The reviewer labels
are entirely human-produced. The source texts
are public-domain and are fetched from Project Gutenberg by identifier.
